\documentclass[11pt]{article}
\input{packages.tex}
\usepackage{ulem}
\DeclareMathOperator{\sign}{sign}
\DeclareMathOperator{\Imm}{Im}
\DeclareMathOperator{\DIF}{DIF}
\DeclareMathOperator{\rg}{rg}
\DeclareMathOperator{\dist}{dist}

% Настройка колонтитулов
\pagestyle{fancy}
\fancyhf{} % Очистить все поля
\fancyhead[L]{А мы тут, знаете, всё \sout{интегралами} плюшками балуемся… }
\fancyhead[R]{\href{https://t.me/mallina_olga }{Автор}}
\fancyfoot[C]{\thepage}               % Справа внизу

\setlength{\headheight}{13.6pt}

\hypersetup{
    colorlinks=true,
    linkcolor=blue,
    citecolor=blue,
    filecolor=blue,
    urlcolor=blue,
}

\begin{document}

\begin{titlepage}
    \centering
    \vspace*{\fill}  % заполняет пространство сверху

    {\Huge Лекция 19. Абсолютная и условная сходимости. Признак Дирихле. Признак Абеля.
    \par}

    \vspace*{\fill}  % заполняет пространство снизу
\end{titlepage}

\setcounter{page}{1}

\section{Лемма об относительной сходимости}
\begin{lemma}
    $\overset{\text{сх}}{\sim}$ определяет отношение эквивалентности на множестве неотрицательных функций на $[a,b)$.
\end{lemma}

\begin{proof}
    1) рефлексивность
    \[ f(x) \overset{\text{сх}}{\sim} f(x), x\to b-0 \]
    \[ m=M=1 \]
    2) симметричность
    \[ f(x) \overset{\text{сх}}{\sim} g(x), \ x\to b-0 \]
    \[\exists m>0, M>0 \text{ т.ч. } mg(x)\le f(x) \le Mg(x) \ \forall x\ge b' \]
    \[ \Rightarrow \frac{1}{M} f(x) \le g(x) \le \frac{1}{m} f(x) \]
    3) транзитивность
    \[ f_1(x) \overset{\text{сх}}{\sim} f_2(x), \ f_2(x) \overset{\text{сх}}{\sim} f_3(x) , x\to b-0 \]
    \[ \Rightarrow f_1(x) \overset{\text{сх}}{\sim} f_3(x), x\to b-0 \]
    \[ f_1(x) \overset{\text{сх}}{\sim} f_2(x) \Leftrightarrow \exists m_1>0,M_1>0 \text{ и } b_1\in (a,b)  \text{ такие, что } m_1f_2(x) \le f_1(x) \le M_1f_2(x) \ \forall x\in [b_1,b) \]
    \[ f_2(x) \overset{\text{сх}}{\sim} f_3(x) \Leftrightarrow \exists m_2>0,M_2>0 \text{ и } b_2\in (a,b)  \text{ такие, что } m_2f_3(x) \le f_2(x) \le M_2f_3(x) \ \forall x\in [b_2,b) \]
    Пусть $b':= max \{b_1,b_2\}$
    \[ m_1m_2f_3(x)\le f_1(x) \le M_1M_2f_3(x) \ \forall x\in [b',b) \]
    \[ \Rightarrow f_1(x)\overset{\text{сх}}{\sim} f_3(x), x\to b-0 \]
\end{proof}

\section{Лемма о неотрицательных функциях}

\begin{lemma}
    Пусть $f_i, g_i \ i=1,2,3$ --- неотрицательные функции на $[a,b)$. Пусть 
    \[ f_i(x) \overset{\text{сх}}{\sim} g_i(x), x\to b-0 \]
    Если $f_3(x)>0$ и $g_3(x)>0 \ \forall x\in [a,b)$, то 
    \[ \frac{f_1(x)f_2(x)}{f_3(x)} \overset{\text{сх}}{\sim} \frac{g_1(x)g_2(x)}{g_3(x)}, \  x\to b-0 \]
\end{lemma}

\begin{proof}
    \[ f_i(x) \overset{\text{сх}}{\sim} g_i(x) \Leftrightarrow \exists m_i>0, M_i>0, b_i\in (a,b) \text{ такие, что } m_ig_i(x)\le f_i(x)\le M_ig_i(x) \ \forall x\in (b_i,b) \]
    Пусть $b':= max \{b_1,b_2,b_3 \}$.
    Тогда
    \[ \frac{m_1m_2g_1(x)g_2(x)}{M_3g_3(x)} \le  \frac{f_1(x)f_2(x)}{f_3(x)} \le \frac{M_1M_2g_1(x)g_2(x)}{m_3g_3(x)} \forall x\in [b',b)\]
    Пусть \[ m= \frac{m_1m_2}{M_3} , M= \frac{M_1M_2}{m_3} \]
    $\Rightarrow$ ч.т.д.
\end{proof}

\section{Абсолютная и условная сходимости несобственного интеграла}

\begin{definition}
    Пусть $-\infty<a<b\le +\infty, f\in \mathcal{R}([a,b']) \ \forall b'\in(a,b) $.
    Будем говорить, что несобственный интеграл 
    \[ \int_{a}^{\to b} f(x)\,dx  \text{ сходится абсолютно, если }  \int_a^{\to b} |f(x)|\,dx \text{ --- сходится.} \]
    Если же 
    \[ \int_{a}^{\to b} f(x)\, dx \text{ --- сходится и при этом  } \int_a^{\to b} |f(x)|\,dx \text{ --- расходится, то говорим, что } \]
    \[ \int_{a}^{\to b} f(x)\,dx \text{ сходится условно. } \]
\end{definition}

\begin{remark}
    Заведём обозначение: $f\in A([a,b)) \Leftrightarrow \forall b'\in (a,b)\ f\in \mathcal{R}([a,b'])$\\
    при этом считаем, что $-\infty<a<b\le +\infty$
\end{remark}

\begin{theorem}
    Пусть $f\in A([a,b))$ . Тогда 
    \[ \int_{a}^{\to b} f(x)\,dx \text{ --- сходится абсолютно } \Rightarrow \int_{a}^{\to b} f(x)\,dx\text{ --- сходится } \]
    Обратное, конечно, неверно. 
\end{theorem}

\begin{proof}
    Ключевое наблюдение: так как $f\in A([a,b)) \Rightarrow |f|\in A([a,b))$ . Кроме того, справедливо неравенство:
    \[ \left| \int_{b'}^{b''} f(x)\,dx \right| \le \int_{b'}^{b''} |f(x)|\,dx \quad \forall b',b''\in [a,b) \tag{*}\label{eq:star} \]
    \[ \int_{a}^{\to b} f(x)\, dx \text{ --- сходится абсолютно} \Leftrightarrow \int_{a}^{\to b} |f(x)|\,dx \text{ --- сходится } \]
    Тогда по критерию Коши:
    \[ \Leftrightarrow \forall \varepsilon>0 \ \exists b(\varepsilon)\in (a,b) : \forall b',b''\in (b(\varepsilon),b) \hookrightarrow \int_{b'}^{b''} |f(x)|\,dx < \varepsilon \]
    Тогда в силу \eqref{eq:star} получим:
    \[ \Rightarrow \forall \varepsilon >0  \ \exists b(\varepsilon)\in (a,b): \forall b',b''\in (b(\varepsilon),b) \hookrightarrow \left|\int_{b'}^{b''}f(x)\,dx \right| <\varepsilon \]
    $\Rightarrow$ выполнено условие Коши для несобственного интеграла 
    \[ \int_{a}^{\to b} f(x)\,dx \]
    $\Rightarrow$ в силу критерия Коши 
    \[ \int_{a}^{\to b} f(x)\,dx \text{ --- сходится }  \]
\end{proof}

\begin{theorem}
    (Признак Дирихле для несобственных интегралов) Пусть $f\in C([a,b)), g\in C^1([a,b))$.Пусть\\
    1) Первообразная F от f ограничена на $[a,b)$\\
    \[ (F(x)=\int_{a}^{x}f(t)\,dt \ ,x\in[a,b)) \] 
    2) $\exists \lim\limits_{x\to b-0} g(x) =0 $\\
    3) $g$ нестрого убывает на $[a,b)$, то есть  $g'(x)\le 0 \ \forall x\in [a,b)$\\
    Тогда
    \[ \int_{a}^{\to b}f(x)g(x)\, dx \text{ --- сходится }\]
\end{theorem}

\begin{proof}
    Фиксируем произвольный $b'\in (a,b)$. Проинтегрируем по частям:
    \[ \int_{a}^{b'} f(x)g(x)\, dx = F(x)g(x) \bigg|_a^{b'} - \int_{a}^{b'}F(x)g'(x)\, dx \]
    Так как $g'$ --- непрерывна, то по формуле Ньютона-Лейбница
    \[ \int_{a}^{b'} g'(x)\,dx = g(b') - g(a) \tag{*}\label{eq:star}\]
    В силу условия 2) и \eqref{eq:star} 
    \[ \Rightarrow \int_{a}^{\to b} g'(x)\, dx \text{ --- сходится } \]
    Но при этом $g'(x)\le 0$, а значит
    \[ \int_{a}^{\to b} -g'(x)\, dx = \int_{a}^{\to b} |g'(x)|\,dx \text{ --- сходится } \]
    В силу условия 1) 
    \[ \exists c>0: |F(x)| \le c \ \forall x\in[a,b) \]
    Но тогда по признаку сравнения:
    \[ \int_{a}^{\to b} |F(x)g'(x)|\, dx \text{ --- сходится } \] 
    \[ \Rightarrow \int_{a}^{b'}F(x)g'(x)\, dx \text{ --- абсолютно сходится }\Rightarrow \text{ сходится } \]
    В силу условия 2)
    \[ \exists \lim\limits_{b'\to b-0} F(b')g(b') =0 \]
    \[ \Rightarrow F(x)g(x) \bigg|_a^{b'} = F(b')g(b') - F(a)g(a) \]
    Таким образом, 
    \[ \exists \lim\limits_{b'\to b-0} F(x)g(x) \bigg|_a^{b'} = -F(a)g(a)\]
    \[ \exists \lim\limits_{b'\to b-0} \int_{a}^{b'} F(x)g'(x)\, dx \in \mathbb{R}\]
    \[ \Rightarrow \exists \lim\limits_{b'\to b-0} \int_{a}^{b'} f(x)g(x)\, dx \in \mathbb{R} \]
\end{proof}

\begin{theorem}
    (Признак Абеля для несобственных интегралов) Пусть $f\in C([a,b)), g\in C^1([a,b))$.\\
    Пусть:
    \[ 1) \int_{a}^{\to b} f(x)\, dx \text{ --- сходится } \]
    \[ 2) g \text{ ограничена снизу на } [a,b) \]
    \[ 3) g \text{ нестрого убывает на } [a,b) \text{, то есть } g'(x)\le 0 \ \forall x\in [a,b) \]
    Тогда 
    \[ \int_{a}^{\to b}f(x)g(x)\, dx \text{ --- сходится }\]
\end{theorem}

\begin{proof}
    Так как $g$ ограничена снизу на $[a,b)$ и нестрого убывает, то 
    \[ \exists \lim\limits_{x\to b-0} g(x) = g_0 \in \mathbb{R} \]
    Рассмотрим функцию $\widetilde{g(x)} := g(x) - g_0 \ \forall x\in [a,b) $. Тогда $\widetilde{g}$ удовлетворяет 2) и 3) в признаке Дирихле.
    $\Rightarrow$ по признаку Дирихле:
    \[ \int_{a}^{\to b} f(x)\widetilde{g(x)}\, dx \text{ --- сходится } \]
    \[ \int_{a}^{\to b} f(x)g_0\, dx \text{ --- сходится по условию 1)} \]
    Но 
    \[ f(x)g(x) = f(x)\widetilde{g(x)} + f(x)g_0 \]
    $\Rightarrow$ в силу линейности несобственного интеграла:
    \[ \int_{a}^{\to b} f(x)g(x)\, dx \text{ --- сходится, и более того он равен =} \int_{a}^{\to b} f(x)\widetilde{g(x)}\, dx + \int_{a}^{\to b} g_0f(x)\, dx \]
\end{proof}


\begin{example}
    Рассмотрим
    \[ I_\alpha = \int_{1}^{+\infty} \frac{\sin x}{x^\alpha}\,dx \]
    Задача: исследовать на абсолютную и условную сходимость при всех $\alpha\in \mathbb{R}$.\\
    1) Заметим, что при $\alpha>1$ 
    \[ \left| \frac{\sin x}{x^\alpha} \right| \le \frac{1}{x^\alpha } \] 
    \[ \int_{1}^{+\infty} \frac{dx}{x^\alpha} \text{ --- сходится при } \alpha>1 \]
    \[ \Rightarrow \text{ по признаку сравнения } \int_{1}^{+\infty} \left| \frac{\sin x}{x^\alpha} \right| \,dx \text{ --- сходится } \]
    2) Заметим, что 
    \[ \frac{1}{x^\alpha} \downarrow 0, x\to +\infty \]
    \[ \left| \int_{1}^{x} \sin t\, dt \right| = \left|\cos t \ \bigg|_1^{x}\right| \le 2 \quad \forall x\in \mathbb{R} \]
    \[ \Rightarrow \text{по признаку Дирихле } \int_{1}^{+\infty } \frac{\sin x}{x^\alpha} dx \text{ --- сходится } \]
    3) Покажем, что при $\alpha \le 0$ 
    \[ \int_{1}^{+\infty } \frac{\sin x}{x^\alpha} dx \text{ --- расходится } \]
    Заметим, что 
    \[ \frac{1}{x^\alpha } \ge 1 \text{ при } \alpha\le 0 \ \forall x\in [1,+\infty) \]
    Запишем отрицание к условию Коши:
    \[ \exists \varepsilon>0: \forall b \in (1,+\infty) \ \exists b',b'' \ge b: \left| \int_{b'}^{b''} \frac{\sin x}{x^\alpha}\, dx \right| \ge \varepsilon  \]
    \[ \exists \varepsilon=1: \forall b \in (1,+\infty) \ \exists b'= 2\pi ([b]+1) ,b'' = \pi+2\pi ([b]+1) : \left| \int_{b'}^{b''} \frac{\sin x}{x^\alpha}\, dx \right| \ge \int_{b'}^{b''} \sin x\,dx = 2 >1 \] 
    \[ \Rightarrow \int_{1}^{+\infty} \frac{\sin x}{x^\alpha} \, dx \text{ --- расходится при } \alpha \le 0 \]
    4) $\alpha\in (0,1] $
    \[ \frac{|\sin x|}{x^\alpha } \ge \frac{\sin^2 x}{x^\alpha} = \frac{1-\cos (2x)}{2x^\alpha} \]
    Теперь рассмотрим следующий интеграл:
    \[ \int_{1}^{+\infty} \frac{1-\cos (2x)}{2x^\alpha} \, dx \text{ --- расходится } \] 
    \[ \text{ так как } \int_{1}^{+\infty } \frac{dx}{2x^\alpha} \text{ --- расходится при } \alpha \in (0,1] \]
    \[ \text{ и } \int_{1}^{+\infty} \frac{\cos (2x)}{2x^\alpha}\, dx \text{ --- сходится по признаку Дирихле (рассуждения те же, что и в пункте 1)}  \]
    Итого: 
    \[ \text{при } \alpha \le 0 \text{ --- расходится } \]
    \[ \text{при } \alpha \in (0,1] \text{ --- сходится условно } \]
    \[ \text{при } \alpha > 1 \text{ --- сходится абсолютно } \]
\end{example}

\begin{remark}
    Условие монотонности функции g в признаке Дирихле очень важно и его нельзя убрать.
\end{remark}

\begin{example}
    Рассмотрим функцию $f(x) = \sin x$ (первообразная ограничена) и \\  
    \[ g(x) = \frac{\sin x}{x}\to 0, x\to +\infty ,\text{ но нет монотонности} \] 
    \[ \int_{1}^{+\infty } f(x)g(x)\, dx \text{ --- расходится } \]
\end{example}

\begin{example}
    (Почему нельзя заменять на эквивалентные в несобственных интегралах от знакопеременных функций ) 
    \[ \int_{2}^{+\infty} \frac{\sin x}{\sqrt{x}} (1+ \frac{\sin x}{\sqrt{x}} )\, dx \]
    \[ \frac{1}{\sqrt{2}} \le 1+ \frac{\sin x}{\sqrt{x}} \le 2 \Rightarrow 1+\frac{\sin x}{\sqrt{x}} \overset{\text{сх}}{\sim} 1 \ ,x\to +\infty \]
    \[ \int_{2}^{+\infty} \frac{\sin x}{\sqrt{x}}\, dx \text{ --- сходится } \]
    \[ \int_{2}^{+\infty} \left(\frac{\sin x}{\sqrt{x}} + \frac{\sin^2 x}{x}\right) \, dx \text{ --- расходится, так как } \int_{2}^{+\infty} \frac{\sin^2 x}{x}\, dx \text{ --- расходится} \]
    Итого, замена на эквивалентную может изменить характер сходимости. 
\end{example}

\begin{theorem}
    (Следствие из признака Абеля). Пусть $f\in C([a,b)), g\in C^1([a,b)) $. Пусть g монотонна на $[a,b)$ и $\exists \lim\limits_{x\to b-0} g(x) = g_0 \neq 0$.
    Тогда 
    \[ \int_{a}^{\to b} f(x)\, dx \text{ и } \int_{a}^{\to b} f(x)g(x)\, dx \text{ имеют один и тот же тип сходимости} \]
\end{theorem}

\begin{proof}
    Пусть 
    \[ \int_{a}^{\to b} f(x)\, dx \text{ --- сходится} \Rightarrow \int_{a}^{\to b} f(x)g(x)\, dx \text{ --- сходится по признаку Абеля} \]
    Пусть 
    \[ \int_{a}^{\to b} f(x)g(x)\, dx \text{ --- сходится} \]
    Так как 
    \[ \exists \lim\limits_{x\to b-0} g(x) = g_0 \neq 0 \Rightarrow \exists b'\in (a,b) \text{ такой, что } g(x) \neq 0 \ \forall x\in (b',b) \]  
    \[ \Rightarrow \frac{1}{g(x)} \text{ корректно определено} \]
    Кроме того, 
    \[ \exists \lim\limits_{x\to b-0} \frac{1}{g(x)} = \frac{1}{g_0} \]
    Кроме того, 
    \[ \frac{1}{g(x)} \text{ тоже является монотонной на } (b',b) \]
    $\Rightarrow$ по признаку Абеля 
    \[ \int_{b'}^{\to b} f(x)\, dx = \int_{b'}^{\to b} f(x) g(x) \frac{1}{g(x)} \, dx \text{ --- сходится по признаку Абеля } \] 
    $\Rightarrow$ в силу принципа локализации 
    \[ \int_{a}^{\to b} f(x)\, dx \text{ --- сходится } \]
    Так как 
    \[ |f(x)g(x)| = |f(x)| |g(x)| \] 
    \[ \text{ и } g(x)\to g_0\neq 0, x\to b-0 \] 
    \[ \Rightarrow |g(x)| \overset{\text{сх}}{\sim} 1, x\to b-0 \]
    $\Rightarrow$ в силу второго признака сравнения
    \[ \int_{a}^{\to b} |f(x)g(x)|\, dx \text{ и } \int_{a}^{\to b} |f(x)|\, dx \text{ --- сходятся или расходятся одновременно } \]
    \[ \Leftrightarrow \int_{a}^{\to b} f(x)g(x)\, dx \text{ --- сходится абсолютно } \Leftrightarrow \int_{a}^{\to b} f(x)\, dx \text{ --- сходится абсолютно} \]
\end{proof}
 
\begin{example}
    Исследовать на сходимость и абсолютную сходимость:
    \[ \int_{1}^{+\infty }x^\alpha \ln (1+\frac{1}{x}) \sin x^2 \]
\end{example}

\begin{proof}
    Рассмотрим следующую функцию:
    \[ g(x) = x\ln \left(1+\frac{1}{x}\right)\to 1, x\to +\infty \]
    \[ g'(x) = \ln \left(1+\frac{1}{x}\right) +  \frac{x}{1+\frac{1}{x}} \cdot \frac{-1}{x^2}  = \ln\left(1+\frac{1}{x}\right) - \frac{1}{1+x}  \] 
    Разложим по Тейлору:
    \[ = \frac{1}{x} - \frac{1}{1+x} - \frac{1}{2x^2} + o\left(\frac{1}{x^2}\right) \]
    \[ \Rightarrow g'(x) =   \frac{1}{x(1+x)} - \frac{1}{2x^2} + o\left(\frac{1}{x^2} \right)\ge \frac{1}{3x^2} \text{ при } x\ge N\in \mathbb{N} \] 
    \[ \Rightarrow g'(x)\ge 0 \ \forall x\ge N \] 
    $\Rightarrow$ в силу следствия из признака Абеля достаточно исследовать
    \[ \int_{1}^{+\infty} x^{\alpha -1 } \sin x^2 \, dx \]
    Пусть $x^2=t, x=\sqrt{t}, 2x\, dx = dt $ 
    \[ \Rightarrow \int_{1}^{+\infty} x^{\alpha -1 } \sin x^2 \, dx \Leftrightarrow \int_{1}^{+\infty} \frac{t^{\frac{\alpha-2}{2}}\sin t}{2} \, dt \]
    Таким образом получаем, что тип сходимости исходного интеграла тот же, что и
    \[ \int_{1}^{+\infty}  t^{\frac{\alpha-2}{2}}\sin t \, dt \]
    \[ \text{ при } \frac{2-\alpha}{2}\le 0 \text{ --- расходится } \]
    \[ \text{ при } \frac{2-\alpha}{2}\in (0,1] \text{ --- сходится условно } \]
    \[ \text{ при } \frac{2-\alpha}{2}>1 \text{ ---  сходится абсолютно} \]
    $\Rightarrow$ итого: 
    \[ \text{ при } \alpha \ge 2\text{ --- расходится } \]
    \[ \text{ при } 0\le \alpha <2 \text{ --- сходится условно } \]
    \[ \text{ при } \alpha<0  \text{ ---  сходится абсолютно} \] 
\end{proof}







\end{document}